\documentclass[11pt]{article}
\usepackage[margin=1in]{geometry}
\usepackage{amsmath,amssymb,amsthm}
\usepackage{booktabs}
\usepackage{enumitem}
\usepackage[hidelinks]{hyperref}
\setlist{nosep,leftmargin=1.5em}
\newtheorem{theorem}{Theorem}[section]
\newtheorem{lemma}[theorem]{Lemma}
\newtheorem{proposition}[theorem]{Proposition}
\newtheorem{corollary}[theorem]{Corollary}
\theoremstyle{definition}
\newtheorem{remark}[theorem]{Remark}
\newcommand{\PG}{\mathrm{PG}}
\newcommand{\Ga}{\mathrm{Ga}}
\newcommand{\R}{\mathbb{R}}
\newcommand{\Lt}{\widetilde L}
\newcommand{\fl}{\mathrm{fl}}
\newcommand{\Log}{\mathrm{Log}}
\newcommand{\lop}{\mathrm{log1p}}
\newcommand{\emo}{\mathrm{expm1}}
\title{The engine's certified Polya-Gamma evaluator:\\ what it adds to the proof, with proofs}
\author{dbarts}
\date{2026-11-07}
\begin{document}
\maketitle

\section{Scope}

The proof note beside this one (``The accept-or-keep Polya-Gamma step'', cited here as \emph{the note}, its
results by name) specifies the evaluator by example in ball arithmetic at 53 bits (Arb). The engine carries the
same enclosures in binary64 with the rules of the note's appendix ``Rounding'', and departs from the example where
the example would be slow or lose accuracy at large shape. This note states each departure and proves it, and
states the cheaper certificate's settings and what they need (the note's ``A cheaper certificate'', whose cells and
$J_0$ S1 was to bound). Every result here is a bound the engine evaluates; none is an approximation. Throughout,
$u=2^{-53}$, $f=\PG(b,\psi)$, $A=|\psi|/2$, $q=\sum_j\Ga(a_j,\mu_j)$ the proposal (the $K$ rates $r_k$ with shape
$b$ and the remainder's $(\alpha,\beta)$), and on the contour $z=\rho+iy$, $\Delta=z^2-\rho^2=-y^2+2i\rho y$.

\section{Arithmetic}\label{sec:arith}

\begin{itemize}
\item[(A1)] As the note's: IEEE 754 binary64, round to nearest; $+,-,\times,/$, $\sqrt{\cdot}$ and the fused
multiply-add correctly rounded.
\item[(A2$_{\mathrm s}$)] The math library returns $\exp,\emo,\log,\lop,\sinh,\cosh,\tanh$ and $\mathrm{atan2}$
within relative error $u_L=2^{-51}$ (two ulps) at normal results; $\sin$ and $\cos$ within $u_L$ absolutely at every
argument; and $\sin$ on $|x|\le\frac12$ and $\cos$ and $\tan$ on $[0,\frac\pi2-2^{-15}]$ within $u_L$ relative.
This adds to the note's (A2), which assumes $\sin$ and $\cos$ only to an absolute bound, the relative accuracy on
those two intervals: the engine's $\emo$ of a small complex ball needs it for $\sin$ and $\cos$, and the closed
forms for $W<0$ (Section~\ref{sec:lines}) for $\cos$ and $\tan$, whose arguments stay below $\frac\pi2-2^{-11}$.
The cell bound and the exponential use $\cos$ and $\sin$ only absolutely, at arguments that reach about $b|y|$
(millions, deep in the left tail). glibc documents two ulps; Apple's library met it on every sample; an argument
reduction with a 66-bit $\pi$ meets the relative part, as it is exact on those intervals. It is an assumption, not a
theorem; Windows' library is not verified.
\item[(A3)] Contraction. A compiler may fuse $a\times b+c$ into one rounding, within an expression or (GCC's
default in GNU modes on targets with a fused multiply-add) across statements. Every ball rule below charges at
least the error of the unfused sequence, and a fused operation's error, $u|ab+c|$, is at most that; so the ball
rules hold either way. The double-double operations rely on error-free transformations, which a fusion across
statements breaks (a product rounded for one use and fused into a sum for another); the header keeps each
transformation's inputs and its rounded product opaque to the compiler, so they hold either way as well, and the
test suite builds the evaluator's tests with contraction forced across statements (built so without the barrier,
every function's error is near $2^{-53}$). The proposal's rates $r_k=\fl(\fl(2\pi^2(k-\frac12)^2)+\psi^2/2)$ are
formed unfused for the same reason: so they are the same doubles on every build and in the reference script.
\end{itemize}
Underflow: each operation that can underflow loses at most $2^{-1074}$ absolutely, and no later operation
magnifies that loss: the modulus bounds $\|m\|^\pm$ scale their components by $2^{\pm600}$, exactly, when
$|m_r|+|m_i|$ is below $2^{-500}$ or above $2^{500}$, so their squares neither underflow nor overflow, and the
quotient scales an operand outside $[2^{-400},2^{400}]$ (in $|m_r|+|m_i|$) to near 1 by a power of two, its radius
rounded up, before dividing; the radii's final $\mathrm{up}(r)
=r(1+2^{-46})+2^{-1020}$ covers that for up to $2^{50}$ operations, and covers the relative rounding of the
up to about 120 operations that form any one radius. The longer sums are charged $1+(N+4)u$ for the rounding of
$N$ nonnegative terms: the main sum's radius (up to about 1,200 nodes), $q$'s sums over a run of terms (up to $K$)
and over runs, and the node-by-node tail (up to $4\times10^5$ nodes). Overflow: an infinite or NaN midpoint or radius makes the
evaluation undecided, never a wrong enclosure.

\paragraph{Complex balls.} A ball is $(m,r)$, $m\in\mathbb C$, the value within $|v-m|\le r$. With $\|m\|^+$ an
upper bound on $|m|$ (the computed $\sqrt{a^2+b^2}$ times $1+2u$) and $\|m\|^-$ a lower one:
\begin{itemize}
\item \emph{Product}, the textbook formula: the midpoint's error is at most $\sqrt5\,u|m_am_b|$ unfused
\cite{bpz} and $2u|m_am_b|$ with one product fused \cite{jklm}; charged $2.25u\|m_a\|^+\|m_b\|^+$, plus the
propagation $\|m_a\|r_b+\|m_b\|r_a+r_ar_b$.
\item \emph{Quotient} $m_a\bar m_b/|m_b|^2$: the numerator's product as above ($\sqrt5u$), the denominator's
relative error at most $2u+u^2$, and the final divisions $u$ per component: at most $(\sqrt5+3)u+O(u^2)\le5.5u$
times $|m_a|/|m_b|$; propagation $(\|m_a\|r_b+\|m_b\|r_a)/(\|m_b\|^-(\|m_b\|^--r_b))$ for $r_b<\|m_b\|^-$.
\item \emph{Exponential} $e^{m}$, $m=\mu+i\nu$: computed as $e^\mu\cos\nu+ie^\mu\sin\nu$: $e^\mu$ within $u_L$
relative moves the midpoint by $u_Le^\mu$, $\cos\nu$ and $\sin\nu$ each within $u_L$ absolutely by at most
$\sqrt2u_Le^\mu$, and the two products by $ue^\mu(1+O(u))$; so the midpoint is within $((1+\sqrt2)u_L+u)e^\mu
(1+O(u))$; propagation $|e^{m+d}-e^m|\le e^\mu\,\emo(|d|)$. Charged $e^\mu(1+u_L)\,[\emo(r)(1+u_L)+2.5u_L+u]$.
\item \emph{$\emo$} for $|\mathrm{Re}|+|\mathrm{Im}|<\frac12$: $\emo(m)=\emo(\mu)\cos\nu-2\sin^2(\nu/2)+i\,e^\mu
\sin\nu$, cancellation-free for small $|m|$; $|\nu|<\frac12$, so $\sin\nu$, $\sin(\nu/2)$ and $\cos|\nu|$ are
within $u_L$ relative under (A2$_{\mathrm s}$), and each of the three terms within $(2u_L+2u)$ of itself, so the
midpoint
within $(2u_L+2u)(|t_1|+t_2+|\mathrm{Im}|)$; propagation $|e^{m}||\emo(d)|\le e^\mu(1+u_L)\emo(r)$. Beyond
$\frac12$, $e^m-1$ from the exponential rule.
\item \emph{$\lop$} $\Log(1+m)$, principal: real part $\frac12\lop(m_r(2+m_r)+m_i^2)$, the argument of $\lop$ within
$3u(|m_r||2+m_r|+m_i^2)$; imaginary part $\mathrm{atan2}(m_i,1+m_r)$ within $u_L$ relative plus the propagated
error of $1+m_r$, at most $2u|m_i|/|1+m|$ (the derivative of $\arg$ along the real axis is $|m_i|/|1+m|^2$ and
$1+m_r$ carries $u|1+m_r|$). Propagation: $|\Log(1+m+d)-\Log(1+m)|\le-\log(1-e)\le e/(1-e)$, $e=r/|1+m|^-$, for
$e<1$, on a ball that misses the cut (the note's ``Branch cuts'').
\item \emph{Square root}, principal, by $p=\sqrt{(|m|+m_r)/2}$, $q=m_i/(2p)$ ($m_r\ge0$; the mirror form for
$m_r<0$): each part within $3.5u(|p|+|q|)$; propagation $r/(\sqrt2\sqrt{|m|^--r})$, the derivative $1/(2\sqrt w)$
bounded on the ball with $|w|\ge|m|-r$ and the factor $\sqrt2$ the bound $|\sqrt a-\sqrt b|\le|a-b|/(\sqrt{|a|}+
\sqrt{|b|})\cdot\sqrt2$ for $a,b$ in one closed half plane.
\end{itemize}
These are the note's rules made explicit at binary64; the constants are the ones in the header.

\section{The target's node in the lean form}\label{sec:lean}

The note's Proposition ``Cancellation-free node'' writes $\log(\Lt(z)/\Lt(\rho))=-b\mathcal D$ with $\mathcal D=
\log\cosh w-\log\cosh w_0$ on the continuous branch, $W_0=A^2+t_0/2$, $W=W_0+\Delta$, $w=\sqrt W$, and its
$\lop(\emo(\delta)(1-e^{-(w+w_0)})/(1+e^{-2w_0}))$, $\delta=\Delta/(w+w_0)$. The node exponent is
$E=2x\Delta-b\mathcal D$. At large $b$ both $2x\Delta$ and $b\mathcal D$ are large and nearly cancel at small $y$ (the
saddle equation), so their separate roundings set the floor. When $w_0\ge1$ the engine uses:

\begin{proposition}[Lean form]\label{prop:lean}
Let $W_0>0$, $w_0=\sqrt{W_0}$, and $y\ne0$. With $\delta=w-w_0=\Delta/(w+w_0)$, $\sigma_0=1/(1+e^{2w_0})$ and
$c_1=4xw_0-b$,
\[
E=2x\delta^2+c_1\delta-b\,\lop\bigl(\sigma_0\,\emo(-2\delta)\bigr),
\]
with the principal $\lop$, which is the continuous branch along the contour.
\end{proposition}
\begin{proof}
$\log\cosh w=w-\log2+\Log(1+e^{-2w})$ is a logarithm of $\cosh w$ for every $w$ with $1+e^{-2w}\notin(-\infty,0]$.
Principal $w$ has $\mathrm{Re}\,w\ge0$, so $|e^{-2w}|\le1$ and $1+e^{-2w}$ lies in the closed disk of radius 1 about
1, meeting $(-\infty,0]$ only at 0, i.e.\ at $e^{-2w}=-1$, $w\in i\pi(\mathbb Z+\frac12)$, $W\le0$ real. On the
contour $\mathrm{Im}\,W=2\rho y\ne0$ for $y\ne0$, and at $y=0$, $W=W_0>0$; so $\Log(1+e^{-2w})$ is continuous along
the contour and real at $y=0$, and $w-\log2+\Log(1+e^{-2w})$ is the branch of $\log\cosh w$ continuous from the real
value at $y=0$ ($w$ itself is continuous: $W$ never meets the principal root's cut, which needs $W\le0$ real). Hence
$\mathcal D=\delta+\Log(1+e^{-2w})-\log(1+e^{-2w_0})$, and since $1+e^{-2w_0}>0$,
$\Log(1+e^{-2w})-\log(1+e^{-2w_0})=\Log\bigl((1+e^{-2w})/(1+e^{-2w_0})\bigr)$ with no change of branch. The quotient
is $1+(e^{-2w}-e^{-2w_0})/(1+e^{-2w_0})=1+\sigma_0(e^{-2\delta}-1)$, as $e^{-2w_0}/(1+e^{-2w_0})=\sigma_0$. Finally
$2x\Delta=2x\delta(w+w_0)=2x\delta(2w_0+\delta)$, so $E=2x\delta^2+4xw_0\delta-b\delta-b\,\lop(\sigma_0\emo(-2\delta))$.
\end{proof}
At the saddle $4xw_0=b\tanh w_0$ (the note's saddle equation), so $c_1=-2b\sigma_0$: the first-order term that
cancelled is gone, and $c_1$ is formed by one fused multiply-add from exact doubles ($4x$ is exact), within
$u|c_1|$. The lead uses $\delta$ from the quotient rule, $\sigma_0$ from $e^{-2w_0}$ ($u_L$, a sum and a quotient,
$(u_L+3u)\sigma_0$ plus $2^{-1074}$ when $e^{-2w_0}$ underflows) and the rules of Section~\ref{sec:arith}.
$\mathrm{Re}\,w\ge0$ gives $|\sigma_0\emo(-2\delta)|<1$ only where $|\emo(-2\delta)|<1+e^{2w_0}$; the $\lop$ rule
needs only that the ball misses the cut, which it checks (an undecided node otherwise). For $w_0<1$ or $W_0\le0$ the
engine uses the note's form unchanged.

\section{The proposal's node in second-order form}\label{sec:qnode}

For $q$, $\log(\Lt(z)/\Lt(\rho))=-\sum_ja_j\Log(1+\zeta_j)$, $\zeta_j=2\Delta/(\mu_j+t_0)$, and at the saddle
$\sum_ja_j/(\mu_j+t_0)=x$.

\begin{proposition}[Second-order form]\label{prop:qnode}
With $c_0=x-\sum_ja_j/(\mu_j+t_0)$,
$E=2x\Delta-\sum_ja_j\Log(1+\zeta_j)=2\Delta c_0-\sum_ja_j\bigl(\Log(1+\zeta_j)-\zeta_j\bigr)$, the principal
$\Log$ being the continuous branch on the contour.
\end{proposition}
\begin{proof}
$2x\Delta=2\Delta c_0+\sum_ja_j\zeta_j$. $1+\zeta_j=(\mu_j+t_0+2\Delta)/(\mu_j+t_0)$ with $\mu_j+t_0>0$, and
$\mathrm{Im}(\mu_j+t_0+2\Delta)=4\rho y$ vanishes only at $y=0$, where the value is positive; so it never meets
$(-\infty,0]$ and the principal $\Log$ is continuous along the contour.
\end{proof}
$c_0$ is computed in double-double (each $a_j/(\mu_j+t_0)$ and the sum). With $\sigma=\sum_ja_j/(\mu_j+t_0)$ and
$n$ terms: each quotient and product is within $17u^2$ of its term, the first addition (to zero) is exact, each of the
other $n-1$ within $3u^2$ of its partial sum, at most $\sigma$ as the terms are positive, and the final subtraction
within $3u^2(\sigma+x)$; so the error is at most $(17+3n)u^2(\sigma+x)$, charged as
$(17+3n)2^{-106}(\sigma+x)(1+n2^{-50})$ with $\sigma$ summed in double, plus $u|c_0|$ for its final rounding to a
double: at the saddle $c_0$ is of the order of the
Newton step's residual, so the first-order term $2\Delta c_0$ carries no large rounding.

\begin{lemma}[$\Log(1+z)-z$ by its series]\label{lem:l1pm}
For $|z|\le m\le\frac14$, $\Log(1+z)-z=-z^2\sum_{k\ge0}(-z)^k/(k+2)$. Truncated after $n$ terms with $m^n\le2^{-54}$,
the remainder is at most $|z|^2m^n/(2(1-m))\le2^{-54}|z|^2/1.5$; Horner's evaluation in complex arithmetic adds at
most $16u|z|^2$ (the engine's loop starts one coefficient high, at $1/(n+1)$, which changes the
$(-z)^n$ coefficient by $1/((n+1)(n+2))$ and leaves the remainder's bound in place); and a ball of radius $r$
about $z$ with $m=|z|+r\le\frac14$ moves the value by at most $r\,m/(1-m)$.
\end{lemma}
\begin{proof}
The series is the Taylor series of $\Log(1+z)-z$, convergent for $|z|<1$. Its tail after $n$ terms is at most
$|z|^2\sum_{k\ge n}m^k/(k+2)\le|z|^2m^n/(2(1-m))$. Horner: each step $p\leftarrow1/(k+2)-zp$ with $|zp|\le m|p|$
and $|p|\le\sum_km^k/(k+2)\le0.603$ at $m=\frac14$; each step's errors (the rounding of $1/(k+2)$, at most $\frac u2$;
the complex product, $\sqrt5um|p|$; the subtraction, $u|p|$) are at most $1.44u$ and are damped by $m$ at each later
step, so $p$ is within $1.44u/(1-m)\le1.92u$; the final product by $z^2$ (two complex products) adds
$2\sqrt5u|z|^2|p|\le2.7u|z|^2$. Together below $5u|z|^2$, charged as $16u|z|^2$. The derivative
of $\Log(1+z)-z$ is $-z/(1+z)$, at most $m/(1-m)$ on the disk.
\end{proof}

\section{The line bounds}\label{sec:lines}

\paragraph{The aggregated majorant with a smaller $J_0$.} The note's Lemma ``Aggregated majorant on a line''
assumes $J_0\supseteq\{j:A_j<4u^2\}$ (here $u$ is the line, not the roundoff). The engine takes $J_0=\{j:A_j\le
2u^2\}$ (an initial run: the $A_j$ increase in $j$), widened by the rounding of $A_j$, and $\kappa$ the largest
upper bound of $k_j=2(A_j-2u^2)/A_j^2$ over $J_1$:

\begin{lemma}[Aggregated majorant, $J_0$ of the nonpositive terms]\label{lem:agg2}
The note's lemma holds for any finite $J_0\supseteq\{j:k_j\le0\}$ with $\kappa=\sup_{J_1}k_j$; and $\kappa\le1/(4u^2)$.
\end{lemma}
\begin{proof}
The note's proof uses of $J_1$ only that $0<k_j\le\kappa$ there: then $-\log(1-k_jY)\le k_jY/(1-k_jY)\le k_jY/(1-
\kappa Y)$ for $\kappa Y<1$. Both hold by the definition of $\kappa$ and of $J_0$. $g(A)=2(A-2u^2)/A^2$ has its
maximum $1/(4u^2)$ at $A=4u^2$. The engine scans $J_1$ from its first term to the first with $A_j>4u^2$ (past which
$g$ decreases), bounding $g$ on each term's enclosure $[A^-,A^+]$ by $2(A^+-2u^2)/(A^-)^2$, and by $1/(4u^2)$ where
the enclosure contains $4u^2$; for $f$ the infinite tail of $J_1$ lies past $4u^2$ and is covered by its first term.
$S=\sum_{J_1}a_jk_j$: for $q$ the sum of the upper bounds; for $f$ the closed form $2b[T_1(W_u)-2u^2T_2(W_u)]$ minus
the $J_0$ terms, with the closed forms' errors charged as stated in the header ($T_1,T_2$ within $2^{-40}$
relative, from $\tanh$, $\cosh$, $\tan$ and $\cos$ under (A2$_{\mathrm s}$), and the series for $|W|\le0.1$). For
$W<0$ the closed forms are used only where $v=\sqrt{-W}$ has $\frac\pi2-v>2^{-11}$, else the line is undecided:
there $\tan v<2^{11}$, so the rounding of $v$ moves $\tan v$ by at most $uv/(\sin v\cos v)$ and $\sec^2v$ by at most
$2uv\tan v\le6434u$ relative, inside $2^{-40}=8192u$ with the functions' and operations' roundings, and $\tan$ and
$\cos$ stay where (A2$_{\mathrm s}$) makes them relative. The saddle's $v_0$ is capped at $\mathtt{0x1.91ffp0}$,
below $\frac\pi2-2^{-11}$, which binds only for $x$ past about $320b$; a frame off the saddle is still valid, and
there it widens the radii: measured at $b$ from $10^{-3}$ to 1 and $x$ up to $10^5b$, at most 1.7 times the saddle
frame's where $\log f>-1000$, at most 4 times where $\log f>-5000$, and much wider or undecided only where $\log f<
-49000$.
\end{proof}
On the contour line ($u=\rho$) every $J_0$ term has $A_j\le2\rho^2$, so $|z^2+\tau_j'^2|^2=A_j^2-2Y(A_j-2\rho^2)+Y^2\ge
A_j^2$ and the term is at most 0. The rounding-widened test for $J_0$ admits $A_j$ up to about $10^{-15}$
relative above $2\rho^2$, so the note's corollary's $Y_0=\max_{J_0}(A_j-2\rho^2)_+$ is tiny but positive; the
corollary's condition on it is met, as $y_1^2\ge1/x$ by the stopping rule below.

\paragraph{The monotone tail's constants.} The note's Corollary ``A monotone tail and an a priori step''(a) needs
$S<4x$ and $Y_1<(1-\sqrt{S/(4x)})/\kappa$ with $\delta=2x-\frac S2(1-\kappa Y_1)^{-2}>0$. The engine takes
$Y_1=(1-\sqrt{S/(2x)})/\kappa$ when $S<2x$ (then $\delta=x$ exactly in real arithmetic) and
$Y_1=\frac12(1-\sqrt{S/(4x)})/\kappa$ when $2x\le S<4x$, both times $1-2^{-30}$, recomputes $\delta$ from the rounded
$Y_1$ with its rounding charged downward, and uses the corollary only if $\delta>0$. The main sum is continued until
$y_1=(N+1)h$ is at least $\max(x^{-1/2},y_{\mathrm{req}})$, $y_{\mathrm{req}}=\frac12(-\rho+\sqrt{\rho^2+2/\delta})$
(the corollary's condition), which also puts $y_1$ past $y_W=\frac12(-\rho+\sqrt{\rho^2+1/x})<\frac1{2\sqrt x}$, so
the note's Theorem ``A priori bound'''s monotone tail applies at any $Y'\ge y_1$.

\paragraph{$\log\sinh$ past the range of a double.} The note's Lemma ``Line majorants''(a) bounds $|\cosh w|\ge
\sinh p$; for $p>710$, $\sinh p$ overflows. The engine bounds $\log\sinh p$ instead:

\begin{lemma}\label{lem:logsinh}
For $p>20$, $\log\sinh p=p-\log2+\lop(-e^{-2p})$, and the computed value less $4u(p+1)$ is a lower bound under
(A1)--(A2$_{\mathrm s}$) before the final charge $u_L|\cdot|$ the callers add.
\end{lemma}
\begin{proof}
$\sinh p=\frac12e^p(1-e^{-2p})$. $p-\log2$ (with $\log2$ rounded, error below $2^{-54}$) is rounded once, within
$u(p+1)$; $\lop(-e^{-2p})$ is below $e^{-40}$ in magnitude, its error far below $u$; the sum rounds once more. In all
at most $3u(p+1)$.
\end{proof}
At small $x$ (the far left tail, which the proposal reaches at small $b$) $\rho$ and the edge lines' $p$ exceed 710;
without this, the cell bounds were NaN and the evaluation undecided, which the step counted as a close call.

\paragraph{The far tail of $q$.} The note's monotone tail for the finite family uses each term's minimum over $y\ge
Y$, which at large $K$ does not decay until $y^2$ passes the remainder's $\tau'^2$. The engine bounds the nodes past
$y_0$ on dyadic blocks $[Y_a,2Y_a]$ of $Y=y^2$:

\begin{lemma}[Dyadic far tail]\label{lem:fartail}
On the contour line let $Y_{\mathrm{end}}=2\max_j(A_j-2\rho^2)_+$. (a) For $Y\ge Y_{\mathrm{end}}$,
$|\Lt(z)/\Lt(\rho)|\le1$, so $\sum_{kh\ge y_a}|G(kh)|\le e^{-2xY_a}[(\rho+y_a)+(\rho/y_a+1)/(4xh)]$ for $y_a\ge
y_W$. (b) On a block $[Y_a,Y_b]$, with $k_j=2(A_j-2\rho^2)/A_j^2$: a term with $k_j\le0$ contributes at most 0; a term
with $k_jY_b\le\frac18$ at most $\frac47a_jk_jY$; any other at most its minimum over the block,
$a_j(\log A_j-\frac12\log\min_{[Y_a,Y_b]}|z^2+\tau_j'^2|^2)$. The sum $B(Y)=C+sY$ of these bounds and $-2xY$ is linear,
at most $C+s\,Y_{a\text{ or }b}$ by the sign of $s$; the block holds at most $(\sqrt{Y_b}-\sqrt{Y_a})/h+2$ nodes, each
with $|z|\le\rho+\sqrt{Y_b}$.
\end{lemma}
\begin{proof}
$|z^2+\tau_j'^2|^2=A_j^2-2Y(A_j-2\rho^2)+Y^2\ge A_j^2$ iff $Y\ge2(A_j-2\rho^2)$ (or $A_j\le2\rho^2$), which (a)'s $Y$
satisfies for every $j$; the node sum is then the note's Lemma ``Line majorants''(c). (b) $-\frac12\log(1-c)\le
\frac c{2(1-c)}\le\frac47c$ for $0\le c\le\frac18$, with $c=k_jY\le k_jY_b$ (the note's (b) bound
$|z^2+\tau_j'^2|^2\ge A_j^2(1-k_jY)$); the per-term minimum is the note's (b). A linear function is largest at an
end. $2100$ doublings reach any $Y_{\mathrm{end}}$ from any positive $y_0^2$ in binary64.
\end{proof}
The blocks' bounds are summed in the log domain (an upper bound on $\log\sum e^{v_i}$ from upper bounds $v_i$, with
the rounding of the exponentials and of the sum charged). Where the monotone tail of the aggregated majorant does
not apply ($S\ge4x$ or $\delta\le0$), the engine evaluates dropped nodes one by one, as the note's Theorem allows,
and every 16 nodes tries the far tail at the current node; it stops when the far tail is below $\frac1{16}$ of the
truncation's budget, or after $4\times10^5$ nodes (then undecided).

\section{The step and the uniform}\label{sec:step}

The uniform is an interval $[k2^{-53},(k+1)2^{-53})$, $k=(w_12^{21})\vee\lfloor w_2/2^{11}\rfloor$ from two 32-bit
Mersenne-Twister words, which is the interval containing the uniform $U$ whose binary expansion begins with those 53
bits; the rest of $U$'s bits are never drawn, as the decision does not depend on them. With $\ell^\pm$ the
enclosure of $\log r=\log f(y)+\log q(x)-\log f(x)-\log q(y)$: accept if $\ell^-\ge\min(0,\log^+((k+1)2^{-53}))$,
keep if $\ell^+\le\log^-(k2^{-53})$ ($-\infty$ at $k=0$), where $\log^\pm$ are the library's $\log$ $l$ widened by
$c=\fl((u_L+u+2^{-98})|l|)$ and $2^{-1000}$; otherwise refine. The widening rounds outward: $|\log v|\le|l|/(1-u_L)$
under (A2$_{\mathrm s}$), so $c$ exceeds $l$'s error by at least $(u+2^{-98}-21u^2)|l|$, and the rounding of
$l\pm c$ is at most $u(|l|+c)\le(u+5u^2)|l|$, less than that as $2^{-98}=256u^2$. A proposal for which the K rule finds no $K$ up to $2^{20}$, or whose $b$ or $\psi$
is not finite, is a close call before any draw. The enclosure's half-width is the four radii plus $4u$ times the sum of
the four midpoints' magnitudes (the sum's rounding). A decision taken is the one the exact $U$ and $r$ give:
accept means
$r\ge U^+>U$ or $r\ge1$; keep means $r\le U^-\le U$ (a tie has probability 0). The probability that the last level
leaves it undecided is at most the measure of $U$ whose interval meets the enclosure of $\log r$ widened by each
threshold's slack: at most twice the half-width (the density of $\log U$ is $U\le1$), plus $(2u_L+2u)|\log U|$ to
first order and $2^{-1000}$ on each side ($U|\log U|\le1/e$, so at most $1.84u_L$ for both), plus the grain. That is the note's floor
$2(2r_f+2r_q)+2^{1-53}$, $r_f,r_q$ the largest radii, with $16u(m_f+m_q)+2u_L+2^{-999}$ added, $m_f,m_q$ the
largest $|\log f|$ and $|\log q|$; a close call keeps $x$ (the note's ``Close calls'').

\section{The cheaper certificate's settings and what they need}\label{sec:cheap}

The levels, fixed in the header before any evaluation: level 0 the note's ``L10'' (one edge line each side at
$a_\pm=\rho/2$, 10 to 40 cells per line, no splitting, discretization target $2^{-56}$ of the estimated sum), level 1
two lines each side ($a_+\in\{\rho/2,\rho/5\}$, $a_-\in\{\rho/2,1\}$), 32 cells, two levels of splitting, target
$2^{-60}$; level 2 the full default (four and six lines, 200 to 3000 cells for $f$, 100 to 600 for $q$, three levels
of splitting, target $2^{-64}$). Each level is a complete certificate; the step intersects them. Every edge
line's distance is capped at $4/\sqrt x$: on a line at distance $a$ the majorant grows about as $e^{2xa^2}$ (the
first-order term vanishes at the saddle), so $h\approx2\pi a/(2xa^2+L)$ for a target of $e^{-L}$, largest near
$a=\sqrt{L/(2x)}$, about $4.7/\sqrt x$ at $L=45$. The cap binds only deep in the left tail, where $\rho\approx
b/(4x)$ is large; there $a=\rho/2$ gave $h$ near $10^{-4}$ and more nodes than the node limit. The note's theorem
holds for any lines inside the strip, so the cap changes cost, not validity.

\paragraph{Cells.} A line's cell count is $\min(c_{\max},\max(c_{\min},10Y\sqrt x))$, with $Y$ the line's truncation
point; so level 0 uses at most 40 cells per line and 80 per density, whatever $(b,\psi,x)$. What the count does not
guarantee is that level 0 reaches its target; where it does not the step goes on to level 1 and 2, which costs time
and never correctness.

\paragraph{$|J_0|$.} For $f$ on the line $\mathrm{Re}\,z=u$, $A_j=\pi^2(j-\frac12)^2+W_u$ with $W_u=W_0+u^2-\rho^2$,
so
\[
|J_0|=\#\{j\ge1:\pi^2(j-\tfrac12)^2\le u^2+\rho^2-W_0\}=\Bigl\lfloor\frac{\sqrt{(u^2+\rho^2-W_0)_+}}\pi+\frac12
\Bigr\rfloor
\]
up to the rounding of $A_j$. With $\rho^2=\max(W_0,(W_0+\pi^2/4)/2)$ (rounded down to 24 bits): on the contour line
$|J_0|=\lfloor\sqrt{2\rho^2-W_0}/\pi+\frac12\rfloor$, which is 1 for $W_0\le\pi^2/4$ and about $w_0/\pi$ above; on the
edge lines $u=\rho/2$ and $u=3\rho/2$ about $w_0/(2\pi)$ and $3w_0/(2\pi)$. At the saddle $w_0$ solves $\tanh(w_0)/w_0=
4x/b$ after the tilt, $W_0=A^2+t_0/2$: near the bulk of $\PG(b,\psi)$, $w_0\approx A$, so $|J_0|\le\frac3{4\pi}|\psi|+1$
($\le37$ at $|\psi|=150$); in the left tail $w_0\approx b/(4x)$ grows without bound, and the engine caps the scan at
$10^5$ terms, past which the line's aggregated bound is not used (the per-cell direct bound of the note's ``Line
majorants'' remains, and costs $O(1)$ per cell). For $q$, $|J_0|$ is the number of rates with
$\frac12(\mu_j+t_0)\le\rho^2+u^2$, which near the bulk is the few smallest rates and in the left tail ($t_0$ large)
all $K+1$ terms, the bound then costing $O(K)$ per cell as the per-term one does.

\paragraph{What was measured.} At the 500 pins of $f$ and 30 of $q$, every finite enclosure at every level
contained its reference and none was undecided after the last level; 200 points in both tails past the pins (down
to $10^{-6}$ of the mean, up to 40 standard deviations above it) were decided; and in the test's 20{,}000 steps over
five $(b,\psi)$ none went past level 0. These are measurements, not proofs. Neither the proof nor soundness needs
level 0 to decide (an undecided level costs the next, and the levels are intersected); that it decides matters only
for the timing, and that it does at every reachable $(b,\psi,x)$ is not proven. The close-call bound needs a weaker
premise, also unproven: that the last level is finite at every reachable point, with radii as measured at the
pins.

\begin{thebibliography}{9}
\bibitem{bpz} R.~P. Brent, C.~Percival and P.~Zimmermann (2007). Error bounds on complex floating-point
multiplication. \emph{Mathematics of Computation} 76, 1469--1481.
\bibitem{jklm} C.-P. Jeannerod, P.~Kornerup, N.~Louvet and J.-M. Muller (2017). Error bounds on complex
floating-point multiplication with an FMA. \emph{Mathematics of Computation} 86, 881--898.
\end{thebibliography}
\end{document}
